% Generated by `taskgen course examples` from dlp.prime_power (seed 5).
% Do not edit: change examples.yaml and regenerate.
\begin{example}[логарифм в подгруппе порядка $3^3$]\label{ex:ph-prime-power}
Элемент $g = 271$ имеет порядок $q^e = 3^{3} = 27$ по модулю $p = 379$ (проверить). Решить задачу дискретного логарифмирования \[271^x \equiv 1\ (\textrm{mod}\ 379),\] находя цифры $x = x_0 + x_{1} \cdot 3 + x_{2} \cdot 3^{2}$, $0 \le x_i < 3$, по одной: каждая цифра "--- логарифм в подгруппе порядка $3$.

\par\noindent\textit{Решение.} $\gamma = g^{q^{e-1}} = 271^{9} \equiv 327 \ne 1$, а $\gamma^{3} = g^{27} \equiv 1$, поэтому $\ord g$ делит $27$, но не делит $9$, т.\,е. $\ord g = 27$, а $\gamma$ имеет порядок $3$. Таблица подгруппы порядка $3$:\[\begin{array}{c|rrr}k & 0 & 1 & 2 \\ \hline \gamma^k & 1 & 327 & 51\end{array}\] Обратный к $271$ по модулю $379$ (расширенный алгоритм Евклида для $(379, 271)$):
\[\begin{array}{r|r|r|r}
r_i & q_i & u_i & v_i \\ \hline
379 &  & 1 & 0 \\
271 &  & 0 & 1 \\
108 & 1 & 1 & -1 \\
55 & 2 & -2 & 3 \\
53 & 1 & 3 & -4 \\
2 & 1 & -5 & 7 \\
1 & 26 & 133 & -186 \\
0 & 2 & -271 & 379 \\
\end{array}\qquad r_i = u_i \cdot 379 + v_i \cdot 271\]

$271^{-1} \equiv -186 \equiv 193\ (\textrm{mod}\ 379)$.
 Если $u_i = h \cdot g^{-(x_0 + \dots + x_{i-1} q^{i-1})}$, то $u_i = g^{x_i q^i + x_{i+1} q^{i+1} + \dots}$ и $w_i = u_i^{q^{e-1-i}} = \gamma^{x_i}$: цифра $x_i$ "--- логарифм $w_i$ по таблице. Начинаем с $u_0 = h$; $u_{i+1} = u_i \cdot c_i^{x_i}$, где $c_i = g^{-q^i}$ ($c_0 = g^{-1}$, $c_{i+1} = c_i^{3}$):\[\begin{array}{c|r|c|r|c|r|r}i & u_i & q^{e-1-i} & w_i & x_i & c_i & u_{i+1} \\ \hline 0 & 1 & 9 & 1 & 0 & 193 & 1 \\ 1 & 1 & 3 & 1 & 0 & 185 & 1 \\ 2 & 1 & 1 & 1 & 0 & 51 & 1 \\ \end{array}\]Последнее $u_{3} = 1$ "--- контроль: все цифры учтены. $x = 0 + 0 \cdot 3 + 0 \cdot 9 = 0$; проверка: $271^{0} \equiv 1\ (\textrm{mod}\ 379)$. (Здесь $h = 1$: все $w_i = 1$, и $x = 0$ сразу.)\[x = 0\ (\textrm{mod}\ 27).\]
\end{example}
